\chapter*{致读者}
\phantomsection
\addcontentsline{toc}{chapter}{致读者}

\begin{enumerate}
\def\labelenumi{\arabic{enumi}.}
\item
  本系列丛书（卷目见第 vii 和 viii 页）从数学的起点开始，并给出完整的证明。原则上，它不要求读者具备特定的数学知识，而只要求对数学推理有一定的熟悉程度，并具备一定的抽象思维能力。尽管如此，它尤其面向那些至少已掌握大学数学课程第一年或第二年内容的人。
\item
  我们选择的阐述方法是公理化的、抽象的，并且通常由一般到特殊。这一选择是由本专著的主要目的所决定的，即要为整个现代数学体系提供坚实的基础。为此，熟悉相当大量的非常一般的概念和原理是必不可少的。此外，证明的要求对主题内容施加了严格固定的顺序。因此，某些考虑的用处不会立即为读者所明了，除非他已经具有相当广博的数学知识；否则，他必须有耐心暂缓判断，直到时机成熟。
\item
  为了减轻这一不利之处，我们经常在正文中插入一些例子，这些例子涉及读者可能已经知道但尚未在本系列中讨论过的事实。这样的例子总是放在两个星号之间：* \ldots{} *。毫无疑问，大多数读者会发现这些例子有助于他们理解正文，并且即使在第一次阅读时也宁愿不跳过它们。当然，从纯粹逻辑的观点来看，省略它们不会有任何不利之处。
\item
  本系列分为若干卷（此处称为“书”）。前六本书有编号，并且一般而言，正文中的每一陈述都只假定前面各卷中已经讨论过的结果为已知。这条规则在每本书内部都成立，但为了阐述的方便，这些书不再按连续的顺序排列。在每本书（或每章）的开头，读者会找到关于其与其他书逻辑关系的精确说明，从而能够确信不存在任何恶性循环。
\item
  每一章的逻辑框架由该章的定义、公理和定理组成。这些是主要需要牢记以备后续使用的部分。较不重要的结果以及那些容易从定理推出的结果被标记为“命题”、“引理”、“推论”、“注记”等。那些在第一次阅读时可以省略的内容用小号字排印。对特别重要的定理的评注偶尔以“附释”的名义出现。
\end{enumerate}

为避免冗长的重复，有时引入仅在某一章或某一章内某一节中有效的记号或缩写是方便的（例如，在只涉及交换环的一章中，“环”一词总是指“交换环”）。这样的约定总是被明确提及，通常是在它们出现的章节的开头。

\begin{enumerate}
\def\labelenumi{\arabic{enumi}.}
\setcounter{enumi}{5}
\tightlist
\item
  正文中的某些段落旨在预先警告读者避免严重错误。这些段落在页边用记号
\end{enumerate}

Z（“危险弯道”）标出。

\begin{enumerate}
\def\labelenumi{\arabic{enumi}.}
\setcounter{enumi}{6}
\item
  习题的设计既使读者能够确信自己已消化正文，又使其注意到那些在正文中没有位置但仍然有趣的结果。最难的习题带有记号 ♦。
\item
  一般而言，我们遵循普遍接受的术语，除非有充分的理由偏离它。
\item
  我们特别努力始终使用严格正确的语言，而不牺牲简洁性。我们尽可能在正文中提请注意语言的滥用，没有这些滥用，任何数学文本都有沦为迂腐，甚至不可读的风险。
\item
  由于原则上正文由理论的教条式阐述组成，它一般不含文献引用。参考文献汇集在历史注记中，通常位于每章末尾。这些注记在适当之处也包含对理论中未解决问题的指示。
\end{enumerate}

每个历史注记之后的参考文献一般只包含在所讨论理论的发展中具有最重要意义的书籍和原始论文。它绝不自诩完备；特别是，仅用于确定优先权的参考文献几乎总是被省略。

至于习题，我们一般认为没有必要指明其来源，因为它们取自许多不同的来源（原始论文、教科书、习题集）。

\begin{enumerate}
\def\labelenumi{\arabic{enumi}.}
\setcounter{enumi}{10}
\tightlist
\item
  对本系列某一部分的引用如下给出：
\end{enumerate}

\begin{enumerate}
\def\labelenumi{\alph{enumi})}
\item
  如果引用的是同一节中提出的定理、公理或定义，则用其编号引用。
\item
  如果它们出现在同一章的另一节中，则在该引用中也引用该节。
\item
  如果它们出现在同一本书的另一章中，则引用章和节。
\item
  如果它们出现在另一本书中，则首先用其书名引用该书。
\end{enumerate}

《结果概要》用字母 R 引用：因此《集合论》R 表示“集合论结果概要”。

